\documentclass[exercice,a4]{thedocument}%
\usepackage{texmaths-tikz}%
\usepackage{pstricks}%
\usepackage{pst-bezier}%
\usepackage{pstricks-add}%
\usepackage{pst-eucl}%
\usepackage{pst-func}%
\usepackage{pst-tree}%
\usepackage{pst-circ}%
\usepackage{tikz}%
\usepackage{pgfplots}%
\usepackage{tkz-tab}%
\usepackage{tikzpagenodes}%
\usepackage{tabularray}%
\usepackage{tabularx}%
\usepackage{multirow}%
\usepackage{multicol}%
\usepackage{scratch3}%
\usepackage{enumitem}%
\usepackage[np]{numprint}%
\usepackage{pas-tableur}%
\usepackage{listings}%
\usepackage{graphicx}%
\usepackage{adjustbox}%
\usepackage{needspace}%
\usepackage{xcolor}%
\usepackage{eurosym}%
\usepackage{diagbox}%
\usepackage{colortbl}%
\usepackage{multido}%
\usepackage{icomma}%
\usepackage[french]{babel}%
\startdatakeys
    \source{26GENMATPO1}
    \BO{2026}
    \cycle{4}
    \competencesDuSocle{CA,RA,CO}
    \objectifsApprentissage{C4A1aa,C4A2ad,C4B2cb,C4D2aa,C4D3ac}
    \motsCles{angle,réciproque de Pythagore,trigonométrie}
\enddatakeys
\definecolor{gris}{gray}{0.5}%
\providecommand{\wideparen}[1]{\overset{\frown}{#1}}%
\begin{document}%
\providecommand{\degres}{^\circ}%
\begin{minipage}[t]{0.45\textwidth}
On donne la figure ci-contre.

BGF est un triangle rectangle en G.

Voici les informations dont on dispose :

AC $= 5$\ \text{cm} \quad BC$ = 4$\ \text{cm}

AB $= 3$\ \text{cm} \quad BG $= 2$\ \text{cm}

$\widehat{\text{BCF}} = 74\degres$
\end{minipage}\hfill
\begin{adjustbox}{max width=0.5\linewidth,max totalheight=255mm}
	\begin{tikzpicture}[>=stealth,baseline={(C.base)}]
		\begin{scope}[shift={(-2:2)},rotate=-2]
			\draw[fill=blue!40](-0.4,0.4)rectangle(0,0);
		\end{scope}
		\draw[fill=blue!40](90:4)--(90:3.6)arc(-90:-16:0.4)node[pos=0.3,below right]{74~°}--cycle;
		\draw[line width=0.3mm](0:0)node[below right]{B}--(180:3)node[left]{A}--(90:4)node[above](C){C}node[midway, above left]{5 cm}--(58:4)node[above right]{F}--(-2:2)node[below right]{G}--(0:0)node[midway, below]{2 cm};
		\draw[line width=0.3mm](0:0)--(90:4)node[midway, sloped]{||}node[pos=0.4, left]{4 cm}(0:0)--(58:4)node[midway, sloped]{||}node[pos=0.4, below right]{4 cm};
		\draw[<->](0,-0.3)--(-3,-0.3)node[midway, below]{3 cm};
	\end{tikzpicture}
\end{adjustbox} \hfill~
%Début PSTricks
%\begin{minipage}{0.5\textwidth}
%\psset{unit=0.9cm,arrowsize=2pt 3}
%\begin{pspicture}(-4.5,-1.5)(3.5,5.8)
%\pspolygon(4;180)(0;0)(5.5;90)%ABC
%\uput[dl](4;180){A}\uput[dr](0;0){B}\uput[ul](5.5;90){C}
%\pspolygon(5.5;90)(5.5;58)(2.9;0)(0;0)(5.5;58)%CFGBF
%\uput[ur](5.5;58){F}\uput[dr](2.9;0){G}
%\psframe(2.9,0)(2.7,0.2)
%\psarc(5.5;90){0.6}{-90}{-16}
%\rput(-0.7,-1.25){\emph{La figure n'est pas en vraie grandeur.}}
%\rput(-2.5,3){5 cm}\rput(-0.4,2.55){4cm}\rput(1.75,2){4 cm} \uput[d](1.45,0){2 cm}
%\def\barre{\psline(-0.05,-0.1)(-0.05,0.1)\psline(0.05,-0.1)(0.05,0.1)}
%\rput{90}(0,2.75){\barre}\rput{54}(1.5,2.4){\barre}\rput(0.6,4.7){$74\degres$}
%\psline[linewidth=0.6pt]{<->}(-4,-0.25)(0,-0.25)\uput[d](-2,-0.25){3 cm}
%\end{pspicture}
%\end{minipage}
%Fin PSTricks


\medskip
Le but de cet exercice est de déterminer si les points A, B et G sont alignés ou non.

\begin{enumerate}
\item On se place dans le triangle CBF.
\begin{enumerate}[label=\textbf{\alph*.}]
\item Justifier que l'angle $\widehat{\text{CFB}}$ mesure $74\degres$.
\item Calculer la mesure de l'angle $\widehat{\text{CBF}}$.
\end{enumerate}
\item Démontrer que le triangle ABC est un triangle rectangle.
\item Dans le triangle rectangle BGF, calculer la mesure de l'angle $\widehat{\text{FBG}}$.
\item Les points A, B et G sont-ils alignés ? Justifier la réponse. \textbf{\textit{Argumenter la réponse en précisant la démarche.}}
\end{enumerate}


\newpage
\medskip
\showthesource
\end{document}
